% !TeX program = xelatex
% ==========================================================================
% Пропоузал семестрового проєкту — курс «Штучний інтелект», УКУ.
% Це ЄДИНИЙ файл, який вам потрібно редагувати. Уся верстка й брендинг
% живуть у ucuproposal.cls — не змінюйте шрифти, поля, кеглі й кольори.
%
% Компіляція: XeLaTeX + Biber (Overleaf робить це автоматично).
% Локально:   latexmk (див. README.md).
% ==========================================================================
\documentclass[ukrainian]{ucuproposal}
% Опції класу (можна комбінувати через кому):
%   ukrainian | english  — мова документа й логотипу (за замовчуванням ukrainian)
%   slogan                — додає гасло УКУ в колонтитул титулки (за замовчуванням вимкнено)

% --------------------------------------------------------------------------
% Метадані титульної сторінки
% --------------------------------------------------------------------------
\projecttitle{Назва вашого проєкту}
\projecttagline{Одне речення про суть проєкту (опційно — можна видалити рядок)}

\teammember{Прізвище Ім'я}
\teammember{Прізвище Ім'я}
% \teammember{Прізвище Ім'я} % третій учасник (опційно)
% \teammember[https://linkedin.com/in/...]{Прізвище Ім'я} % опційне клікабельне посилання на профіль (LinkedIn тощо)

\mentor{} % ПІБ ментора, якщо є; якщо немає — залиште порожнім, рядок не друкується

\track{research} % research | engineering — не друкується на титулці, лише для внутрішньої категоризації

\repo{https://github.com/team/project} % опційно; якщо репозиторію ще немає — залиште порожнім, рядок не друкується

\proposaldate{2026-09-21}

\begin{document}
\maketitlepage

\section{Проблема та мотивація}
% Що саме ви розв'язуєте і чому це важливо й цікаво в контексті ШІ?
% Орієнтовно 1 абзац.
Опишіть проблему, яку розв'язує ваш проєкт, та її значущість.

\section{Огляд джерел та наявних рішень}
% Що вже зроблено у світі за цією темою? Спирайтесь на references.bib,
% цитуйте через \cite{}. Орієнтовно 3-5 джерел.
Стисло опишіть наявні підходи та роботи, на які ви спираєтесь~\cite{example2023,example2022}.

\section{Дані}
% Які набори даних, звідки, ліцензія. Якщо збираєте самі — як саме.
% Якщо дані не потрібні — поясніть чому.
Опишіть джерела даних, їх обсяг і ліцензію.

\section{Підхід і методи}
% Які алгоритми/архітектури плануєте використати, чи спираєтесь на
% існуючі реалізації та як плануєте їх покращити.
Опишіть плановий підхід і методи.

\section{Оцінювання результатів}
% Метрики, baseline, спосіб порівняння, якісні очікування.
Опишіть метрики та baseline, з яким порівнюватиметесь.

\section{План і ризики}
% Етапи до midterm і до фінального захисту, ключові ризики та план Б.
Опишіть план роботи та основні ризики.

% Приклад врізки з ключовими цифрами (опційно, можна видалити):
% \begin{keyfacts}
% \textbf{Ключові цифри:} 3 датасети \textperiodcentered\ 2 baseline-моделі \textperiodcentered\ 4 тижні до midterm
% \end{keyfacts}

\section{Contribution statement}
% Хто за що відповідає, розподіл у відсотках (сума має дорівнювати 100%).
% Формат відсотка обов'язково NN\% (з екранованим \%), інакше збірка зламається.
\begin{contributions}
  \contribution{Прізвище Ім'я}{50\%}{дані, препроцесинг, EDA}
  \contribution{Прізвище Ім'я}{50\%}{моделювання, евалюація, документація}
\end{contributions}

% Примітка: декларація використання AI-інструментів (AI usage) до цього
% пропоузалу не входить — вона живе в README.md вашого репозиторію.

\ucuprintbibliography

\end{document}
